\documentclass[11pt]{article}
\usepackage[a4paper,margin=1.9cm]{geometry}
\usepackage{fontspec}
\setmainfont{Latin Modern Roman}
\usepackage{titlesec}
\usepackage{enumitem}
\usepackage[table]{xcolor}
\usepackage{longtable}
\usepackage{booktabs}
\usepackage{array}
\usepackage{pifont}
\usepackage{amssymb}
\usepackage{tikz}
\usepackage{tcolorbox}
\usepackage{hyperref}
\usepackage{fancyhdr}
\usepackage{parskip}

\definecolor{gxteal}{HTML}{0F766E}
\definecolor{gxcyan}{HTML}{0284C7}
\definecolor{gxslate}{HTML}{334155}
\definecolor{gxamber}{HTML}{B45309}
\definecolor{gxred}{HTML}{B91C1C}
\definecolor{gxgreen}{HTML}{15803D}
\definecolor{gxviolet}{HTML}{6D28D9}

\hypersetup{colorlinks=true, linkcolor=gxteal, urlcolor=gxcyan,
  pdftitle={GlobeX -- PS 26132 Public Dataset Landscape}, pdfauthor={GlobeX Team}}

\titleformat{\section}{\large\bfseries\color{gxteal}}{\thesection}{0.6em}{}
\titleformat{\subsection}{\normalsize\bfseries\color{gxslate}}{\thesubsection}{0.6em}{}
\titlespacing*{\section}{0pt}{1.7em}{0.7em}

\pagestyle{fancy}
\fancyhf{}
\renewcommand{\headrulewidth}{0.4pt}
\fancyhead[L]{\small\color{gxslate}GlobeX --- PS 26132 Dataset Landscape}
\fancyhead[R]{\small\color{gxslate}SIH 2026}
\fancyfoot[C]{\small\thepage}

\newcommand{\yes}{{\color{gxgreen}\ding{108}}}
\newcommand{\pmid}{{\color{gxcyan}\ding{109}}}
\newcommand{\non}{{\color{black!22}\ding{109}}}
\newcommand{\bld}{{\color{gxamber}\ding{53}}}

\newtcolorbox{callout}[1]{colback=gxteal!5, colframe=gxteal!55, boxrule=0.6pt, arc=2pt,
  left=6pt, right=6pt, top=5pt, bottom=5pt, fonttitle=\bfseries\small, coltitle=white,
  colbacktitle=gxteal, title={#1}}
\newtcolorbox{uspbox}[1]{colback=gxviolet!5, colframe=gxviolet!55, boxrule=0.6pt, arc=2pt,
  left=6pt, right=6pt, top=5pt, bottom=5pt, fonttitle=\bfseries\small, coltitle=white,
  colbacktitle=gxviolet, title={#1}}
\newtcolorbox{warnbox}[1]{colback=gxamber!6, colframe=gxamber!65, boxrule=0.6pt, arc=2pt,
  left=6pt, right=6pt, top=5pt, bottom=5pt, fonttitle=\bfseries\small, coltitle=white,
  colbacktitle=gxamber, title={#1}}
\newtcolorbox{redbox}[1]{colback=gxred!6, colframe=gxred!65, boxrule=0.6pt, arc=2pt,
  left=6pt, right=6pt, top=5pt, bottom=5pt, fonttitle=\bfseries\small, coltitle=white,
  colbacktitle=gxred, title={#1}}
\newtcolorbox{codebox}{colback=gxslate!5, colframe=gxslate!40, boxrule=0.5pt, arc=1pt,
  left=6pt, right=6pt, top=4pt, bottom=4pt, fontupper=\ttfamily\footnotesize}

\begin{document}

% ================= TITLE =================
\begin{titlepage}
\centering
\vspace*{1.8cm}
{\Huge\bfseries\color{gxteal} PS 26132 --- Public}\\[0.4em]
{\Huge\bfseries\color{gxteal} Dataset Landscape}\\[1.8em]
{\Large\color{gxslate} Every source, what it actually contains, how to pull it,}\\[0.3em]
{\Large\color{gxslate} and exactly where the public data runs out}\\[3em]
{\normalsize\color{gxslate} Market linkages \& price discovery for farmers --- Govt.\ of Maharashtra (MSInS)}\\[0.3em]
{\normalsize\color{gxslate} Prepared for GlobeX --- 2 September 2026}\\[2.6em]
\begin{tabular}{cccc}
\textbf{\color{gxgreen}\Large \yes} & \textbf{\color{gxcyan}\Large \pmid} & \textbf{\color{gxamber}\Large \bld} & \\
{\small prices \& arrivals} & {\small quality standards} & {\small FPO/buyer \& storage data} & \\
{\small genuinely public} & {\small published, not dataset-ified} & {\small does not exist publicly} & \\
\end{tabular}
\vfill
{\footnotesize\color{gxslate} Internal working document --- not for external distribution as-is}
\end{titlepage}

% ================= WHAT'S NEEDED =================
\section{What this problem statement actually needs data for}
26132 asks for six things, and public-data coverage is wildly uneven across them. This report goes
source-by-source only for what's real --- no dataset is listed here without a link and a description of
what's actually in it.

\renewcommand{\arraystretch}{1.3}
\noindent\begin{longtable}{@{}p{0.5cm}p{4.6cm}p{2.6cm}p{7.3cm}@{}}
\toprule
& \textbf{Data need} & \textbf{Public GoI coverage} & \textbf{Verdict} \\
\midrule
\endhead
1 & Mandi prices (min/modal/max) & \yes\ Strong & Daily, national, years of history, documented API. \\
2 & Arrival volumes & \yes\ Strong & Same source, same access path as prices. \\
3 & Quality/grading standards & \pmid\ Partial & The AGMARK grade specifications themselves are public; there is no dataset of actual graded lots. \\
4 & FPO directory \& buyer credentials & \bld\ Absent & No structured, complete public dataset exists anywhere in government data. \\
5 & Storage/warehouse options & \bld\ Mostly absent & WDRA publishes a registered-warehouse list; no public capacity or pricing data. \\
6 & Weather / yield signals (optional enrichment) & \pmid\ Partial & IMD and Bhuvan/ISRO data exist but need separate integration; not core to the MVP. \\
\bottomrule
\end{longtable}

% ================= PRIMARY SOURCES =================
\clearpage
\section{Primary Government of India sources --- prices \& arrivals}

\subsection*{Agmarknet, via the Open Government Data (OGD) platform}
Agmarknet itself (\path{agmarknet.gov.in}) is the Directorate of Marketing \& Inspection's public
dashboard for mandi prices, but the portal has no clean bulk API of its own --- integrations go through its
official mirror on \path{data.gov.in}. Two resources matter here:

\begin{itemize}[leftmargin=1.4em]
\item \textbf{``Current daily price of various commodities from various markets (Mandi)''} --- state,
district and market-level daily records: commodity, variety, min/modal/max price (Rs./quintal) and arrival
date. This is the workhorse dataset for the price-trend and sale-window features.
\item \textbf{``Variety-wise Daily Market Prices Data of Commodity''} --- the same data broken out at
variety granularity (e.g.\ not just ``Wheat'' but the specific variety traded), useful when the grading
layer needs to match a specific variety's price rather than a commodity average.
\end{itemize}

\begin{codebox}
GET https://api.data.gov.in/resource/9ef84268-d588-465a-a308-a864a43d0070\\
    ?api-key=\{your\_key\}\&format=json\&offset=0\&limit=100\\
    \&filters[state]=Maharashtra\&filters[commodity]=Onion
\end{codebox}

\textbf{Coverage:} national, state-wise and market-wise; several years of daily history. \textbf{Update
frequency:} daily. \textbf{Access:} free API key from a \path{data.gov.in} account (register with a Google
or government email); a shared public demo key exists for prototyping but carries tighter rate limits and
should not be relied on for a live demo. \textbf{Format:} JSON or CSV via the \texttt{format} parameter,
though some older resources ignore the request and always return CSV --- check the response
\texttt{Content-Type} rather than assuming. \textbf{License:} Government Open Data License -- India.

\subsection*{AIKosh --- the IndiaAI Mission's dataset repository}
\path{aikosh.indiaai.gov.in} is a newer GoI platform, run under the IndiaAI Mission (MeitY), whose explicit
purpose is hosting datasets curated for AI training and use --- as distinct from data.gov.in's general
open-data mandate. It carries a packaged version of the variety-wise mandi price dataset. Citing AIKosh
alongside data.gov.in is worth doing deliberately in the pitch: it is the government's own ``data curated for
AI'' platform, which answers a judge's ``is this real AI on real data'' question about as directly as a
citation can.

\begin{warnbox}{Don't over-rely on any single mirror}
All three of Agmarknet, its data.gov.in mirror, and its AIKosh repackaging trace back to the same underlying
DMI/Ministry of Agriculture source. They are not independent datasets --- treat them as one source with three
access points, and validate whichever one you build the ingestion pipeline against.
\end{warnbox}

% ================= MIRRORS =================
\clearpage
\section{Non-official mirrors --- faster to prototype against}

\renewcommand{\arraystretch}{1.35}
\noindent\begin{longtable}{@{}p{3.6cm}p{7.4cm}p{3.6cm}@{}}
\toprule
\textbf{Source} & \textbf{What it actually offers} & \textbf{Best use} \\
\midrule
\endhead
CEDA Agri Market Data (Ashoka University CEDA) --- \path{agmarknet.ceda.ashoka.edu.in} & Repackages the same
DMI source with cleaner state/district-level exports, Census-2011-mapped boundaries, and built-in
timeseries/heatmap visualisation and comparison tools. & Fast exploratory analysis and sanity-checking
before writing the live ingestion pipeline. \\[0.4em]
India Data Portal --- APMC arrivals and prices (Agmarknet) dataset & A cleaned, versioned CKAN-hosted copy
of the arrivals-and-prices data with a documented resource page. & A stable fallback if the live
data.gov.in API misbehaves during development or demo prep. \\[0.4em]
Kaggle mirrors (e.g.\ ``Indian Agricultural Mandi Prices'', ``Daily Wholesale Commodity Prices -- India
Mandis'') & Community-maintained CSV snapshots of the same underlying data, already flattened. & Offline
development and model prototyping before the live API integration is wired up. \\
\bottomrule
\end{longtable}

% ================= ENAM =================
\section{eNAM --- what it is, and what it is not}
\begin{redbox}{eNAM is a dashboard, not a data source}
\path{enam.gov.in} is the National Agriculture Market trading platform, and its price dashboard
(\path{enam.gov.in/web/dashboard/agmarknet}) simply re-displays the same Agmarknet/DMI data described above
--- it is not an independent feed. eNAM does not expose transaction-level trading data or an open API to
third parties. Any pitch claiming ``eNAM integration'' as a built feature will not survive a judge asking to
see the integration; the honest framing is a future partnership/onboarding target, not a current data
source.
\end{redbox}

% ================= GAPS =================
\clearpage
\section{The real gaps --- where public data simply does not exist}

\subsection*{FPO and verified-buyer directories}
There is no government dataset --- structured or otherwise --- listing farmer producer organisations with
verifiable attributes (registration status, commodity focus, aggregation capacity) or buyer credentials
(procurement history, payment reliability) in a form usable for matching. NABARD publishes some FPO
programme information and individual state cooperation departments maintain scattered registers, but none
of it is complete, structured, or centrally accessible. \textbf{This has to be self-built}: a small,
clearly-labelled pilot seed dataset (real-structure, synthetic-content) is the honest path, not a claim of
sourced data.

\subsection*{Quality grading}
AGMARK grade specifications (the standards themselves --- moisture content, size, foreign matter thresholds,
etc.\ per commodity) are published and public. What does not exist is a dataset of \emph{actual graded
lots} --- no public record of ``this specific batch was graded X.'' Document/OCR verification against the
published standard is honest engineering; claiming a trained grading model would require a labelled dataset
that has to be built from scratch, most plausibly by photographing and grading a small sample set during
prototyping.

\subsection*{Storage and logistics}
The Warehousing Development \& Regulatory Authority (WDRA) publishes a list of registered warehouses, which
gives location and registration status, but not live capacity, current occupancy, or storage pricing.
Transport/freight cost data is likewise not published in a usable open form. The existing GlobeX
ETA/landed-cost estimator can supply distance and time; storage-cost and live-capacity inputs will need to
be assumption-based or manually sourced for the demo, not pulled from a live feed.

% ================= DATA QUALITY =================
\clearpage
\section{Data quality risks --- and the cleaning checklist that avoids them}

\renewcommand{\arraystretch}{1.3}
\noindent\begin{longtable}{@{}p{4.4cm}p{5.4cm}p{4.8cm}@{}}
\toprule
\textbf{Known issue} & \textbf{Why it happens} & \textbf{Cleaning step} \\
\midrule
\endhead
Inconsistent commodity/variety naming across states & Each state market board enters its own free-text
labels into the same national schema & Build a canonical commodity/variety lookup table; fuzzy-match
incoming labels against it before ingestion. \\[0.4em]
Missing days / gaps in the series & Not every market reports every day, and reporting reliability varies by
state & Define an explicit gap-fill policy (carry-forward with a flag, or interpolate) rather than letting
gaps silently become zeros or nulls the forecaster misreads. \\[0.4em]
Non-standard units & Prices are per quintal in most records, but not guaranteed consistent across every
state/market entry & Normalise to one unit at ingestion and validate the conversion, don't assume. \\[0.4em]
API pagination edge case & The API can return an empty \texttt{records} array before the reported
\texttt{total} count is reached & Stop pagination on the first empty response rather than trusting the
total-count arithmetic. \\[0.4em]
Format parameter sometimes ignored & Some older data.gov.in resources always return CSV regardless of the
requested \texttt{format} & Check the response \texttt{Content-Type} header and branch the parser
accordingly, rather than assuming JSON. \\[0.4em]
Field-name case sensitivity in filters & \texttt{filters[State]} works, \texttt{filters[state]} silently
returns nothing on some resources & Confirm exact field casing per resource before relying on server-side
filtering; don't assume a naming convention holds across datasets. \\
\bottomrule
\end{longtable}

\begin{warnbox}{Budget real time for this, not model time}
None of the issues above are modeling problems. They are exactly the unglamorous data-engineering work that
decides whether a live demo works in front of judges. Treat the ingestion/cleaning layer as its own
work item with its own time budget, not an afterthought bolted onto the forecasting task.
\end{warnbox}

% ================= FIELD MAPPING =================
\clearpage
\section{Recommended field mapping into GlobeX's existing schema}
The ingestion pipeline's job is to make mandi records look like the trade-panel records the forecaster and
ranking engine already expect --- same shape, different domain.

\renewcommand{\arraystretch}{1.3}
\noindent\begin{longtable}{@{}p{4.6cm}p{4.6cm}p{5.4cm}@{}}
\toprule
\textbf{Agmarknet/OGD field} & \textbf{Maps to (GlobeX)} & \textbf{Notes} \\
\midrule
\endhead
State, District, Market & Corridor origin (region/market level) & A ``corridor'' becomes a market; no
schema change needed, only relabeling. \\[0.3em]
Commodity, Variety & Product/commodity identifier & Requires the canonical lookup table described above
before it can join cleanly to other tables. \\[0.3em]
Arrival Date & Transaction date / time-series index & Same role it plays for trade data --- the axis the
forecaster runs against. \\[0.3em]
Min/Modal/Max Price & Price features (point + spread) & The min--max spread is a genuinely useful extra
feature the trade-corridor data didn't have --- it's a built-in volatility signal, worth keeping rather than
collapsing to modal price alone. \\[0.3em]
Arrival Quantity & Volume feature & Feeds both the forecaster (demand proxy) and the ranking engine
(glut-risk dimension). \\
\bottomrule
\end{longtable}

\section*{Bottom line}
The one dataset this problem statement structurally cannot exist without --- mandi prices and arrivals --- is
genuinely, solidly public: documented, API-accessible, updated daily, with GoI-run mirrors (AIKosh) as well
as faster third-party ones (CEDA, India Data Portal) for prototyping. Everything else the PS asks for ---
verified buyers, FPO directories, quality-graded lots, storage capacity --- has no public dataset behind it
at all, and needs to be honestly scoped as pilot/seed data or a stated roadmap item rather than implied as
already sourced. That split is exactly what should drive the build order: ingestion and forecasting first,
because the data to do them well actually exists; matching and grading second, built openly on small,
labelled seed data rather than quietly on synthetic data dressed up as real.

\end{document}
